\section{2025 Algebra \& Number Theory}

\subsection*{Exam}
\begin{exercise}
Let $F$ be a field of characteristic $\neq 2$. Show that $F$ has a Galois extension of group $\mathbb{Z}/2\mathbb{Z} \times \mathbb{Z}/2\mathbb{Z}$ iff $[F^*: F^{*2}] > 2$. Show these are splitting fields of $X^4 + aX^2 + b$ with $b \in F^{*2}$.
\end{exercise}

\begin{solution}
\textbf{Motivation:} This exercise connects Kummer theory (which classifies abelian extensions) with the concrete form of biquadratic polynomials.

By Kummer theory, $(\mathbb{Z}/2\mathbb{Z})^n$ extensions of $F$ correspond to $n$-dimensional $\mathbb{F}_2$-subspaces of $F^*/F^{*2}$.
A 2-dimensional subspace exists if and only if $\dim_{\mathbb{F}_2}(F^*/F^{*2}) \geq 2$, which means the index $[F^*:F^{*2}]$ is at least $2^2=4$. The condition $>2$ for a power of 2 is equivalent to $\geq 4$.

For the polynomial $f(X) = X^4 + aX^2 + b$, let $\alpha^2$ and $\beta^2$ be the roots of the quadratic $Z^2 + aZ + b = 0$. The roots of $f$ are $\pm \alpha, \pm \beta$.
The splitting field is $F(\alpha, \beta)$. Since $\alpha\beta = \sqrt{b}$, we have $F(\alpha, \beta) = F(\alpha, \sqrt{b}/\alpha) = F(\alpha, \sqrt{b})$.
For this to be a $V_4$ extension, we need $F(\alpha)/F$ to be quadratic and $\sqrt{b}$ to be in $F$ but not a square of a root of the quadratic? No.
Actually, the Galois group of a biquadratic $X^4+aX^2+b$ is $V_4$ iff $b$ is a square in $F$ and $f$ is irreducible. If $b=k^2$, the roots are $\pm \sqrt{\frac{-a \pm \sqrt{a^2-4k^2}}{2}}$. The extension is $F(\sqrt{u}, \sqrt{v})$ for some $u, v \in F$.
\end{solution}

\begin{exercise}
Let $p$ be prime number, $e, f \in \mathbb{Z}_{\geq 1}$. Let $G$ be a finite group of order $n := ef$ generated by two elements $\sigma$ and $\tau$, satisfying the relations
\[
\sigma^f = 1, \quad \tau^e = 1, \quad \text{and} \quad \sigma\tau\sigma^{-1} = \tau^p.
\]

\begin{enumerate}
\item[(1)] Show that such a group $G$ exists if and only if $p^f \equiv 1 \pmod{e}$. If the latter condition is satisfied, there exists, up to isomorphisms, a unique finite group $G$ as described above.

\item[(2)] Let $\mathbb{F}_{p^f}$ denote the finite field with $p^f$ elements, and $\eta \in \mathbb{F}_{p^f}$ a primitive $e$-th root of unity. Let $G$ act on $A := \mathbb{F}_{p^f}$ via
\[
\sigma(a) = a^p, \quad \tau(a) = \eta \cdot a.
\]
Show that the following three assertions are equivalent:
\begin{enumerate}
\item[(a)] The $\mathbb{F}_p$-representation $A$ of $G$ is absolutely irreducible in the sense that the $\mathbb{C}$-representation $\mathbb{C} \otimes_{\mathbb{F}_p} A$ of $G$ is irreducible, where $\mathbb{C}$ is any algebraically closed field containing $\mathbb{F}_p$;
\item[(b)] the $\mathbb{F}_p$-representation $A$ of $G$ is irreducible;
\item[(c)] $p$ is of order $f$ in $(\mathbb{Z}/e\mathbb{Z})^\times$.
\end{enumerate}
\end{enumerate}
\end{exercise}

\begin{solution}
\textbf{Motivation:} This problem links group theory (semidirect products) with representation theory over finite fields.

(1) The condition $p^f \equiv 1 \pmod e$ ensures that $\sigma$ acts as a valid automorphism of $\langle \tau \rangle$ and that its $f$-th power is identity. This is the condition for the existence of the semidirect product $C_e \rtimes C_f$.

(2) The representation $A$ has $\tau$ acting as multiplication by $\eta$ (a primitive $e$-th root) and $\sigma$ as Frobenius $x \mapsto x^p$.
- If $p$ has order $f \pmod e$, then $\mathbb{F}_p(\eta) = \mathbb{F}_{p^f}$. Any $G$-invariant subspace is an $\mathbb{F}_p(\eta)$-module, hence a vector space over $\mathbb{F}_{p^f}$. Only $\{0\}$ and $A$ exist.
- Absolute irreducibility over $\mathbb{F}_p$ means the only endomorphisms are scalars. Commuting with $\tau$ makes the map $\mathbb{F}_{p^f}$-linear. Commuting with $\sigma$ restricts scalars to $\mathbb{F}_p$.
\end{solution}

\begin{exercise}
(1) $\frac{1}{T-\alpha_i}$ basis of $R/(T^n)$.
(2) Valuation set density $u_n/n \geq 1/[K:k]$.
\end{exercise}

\begin{solution}
(1) These are $n$ elements in an $n$-dimensional space. Their linear independence follows from the uniqueness of partial fraction decompositions. If a sum $\sum \frac{c_i}{T-\alpha_i}$ is $0 \pmod{T^n}$, it has a zero of order $n$ at $T=0$, which is impossible for a rational function of this degree.

(2) This is a clever application of the inequality between $k$-dimension and $K$-dimension. The set $\Omega$ of valuations tracks the "jumps" in the Taylor expansion at $\beta$. The density bound $1/d$ reflects that $K$ is at most $d$ times "larger" than $k$.
\end{solution}

\begin{exercise}
Galois descent for unramified $p$-adic extensions.
\end{exercise}

\begin{solution}
\textbf{Motivation:} Galois descent is a fundamental tool in arithmetic geometry. The unramified case is "nice" because it reduces to the theory of finite fields.

(1) The isomorphism $O_E \otimes_{O_F} O_E \cong \prod O_E$ is the definition of a Galois cover in the context of rings of integers. It holds exactly when the extension is unramified (no discriminant in the residue field).

(2-3) This is the statement that "descent is effective". Any module over $O_E$ with a compatible $G$-action comes from a unique module over $O_F$.

(4) Ramification breaks this because the ring $O_E$ becomes "too large" (the different is non-trivial). The length of the modules doesn't match up under the tensor product functor.
\end{solution}

\begin{exercise}
Splitting field of $X^5 - X + 1$.
\end{exercise}

\begin{solution}
\textbf{Motivation:} This is a concrete example of a group $S_5$ appearing as a Galois group, with interesting ramification properties.

(1) Irreducibility is checked mod 2 and 5.
(2) Ramification only at primes dividing the discriminant $D = 19 \times 151$.
(3) The inertia groups are transpositions. Transpositions are odd permutations, so they are not in $A_5$. Since $K = \mathbb{Q}(\sqrt{D})$ is the fixed field of $A_5$, the inertia groups act trivially on the residue field extension of $E/K$. Thus $E/K$ is unramified at finite primes.
(4-5) The Galois group contains a transposition (inertia) and acts transitively on 5 roots, hence $G = S_5$. The unique quadratic subfield is $\mathbb{Q}(\sqrt{D})$.
\end{solution}
